% ============================================================
% ［架空の記入例・提出不可］研究設定，データ，数値および結果はすべて架空である．
% 引用文献は実在するが，本例の研究成果を裏づけるものではない．
% 通常のchapters/01-introduction.texには自動で読み込まれない．
% ============================================================
\chapter{\UsamiIntroductionTitle}

\begin{center}
  \UsamiFictionalExampleNotice
\end{center}

\section{研究背景}

人物姿勢推定は，画像中の肩，肘，膝等の関節位置を求め，人物の動作を骨格として表現する技術である．高解像度表現を維持するHRNet\cite{sun2019hrnet}やTransformerを用いるViTPose\cite{xu2022vitpose}により，明瞭な人物画像に対する推定精度は向上している．一方，監視映像や夜間の行動計測では，照度不足によるコントラスト低下，動きぼけ，人物間の重なりが同時に生じる．このとき，存在しない位置へ関節が集まった骨格や，左右の手足が入れ替わった骨格が出力され，人数計測や動作分類等の後段処理へ誤りが伝播する．

姿勢推定器を対象環境ごとに再学習する方法は有力であるが，十分な低照度画像と関節アノテーションを準備できない場面もある．また，推定器が出力する関節信頼度だけで骨格を選別すると，一部の関節が高信頼度であっても人体として不自然な骨格を採用する場合がある．そこで，既存の推定器を変更せず，出力された骨格の構造的一貫性を利用して誤推定を除外する後処理を検討する．

\section{研究目的}

本研究の目的は，低照度の単眼RGB動画を対象として，関節信頼度と骨格辺長の一貫性から誤推定骨格を判定する方法を構築し，精度と採用率の両面から有効性を評価することである．ここで誤推定の除外とは，関節座標を補正することではなく，後段処理へ渡す骨格を選択することを指す．したがって，採用した骨格の位置精度が高くても，多数の骨格を捨てる方法は実用上十分とは限らない．本研究では，検証系列だけで判定閾値を決め，未使用の評価系列においてPCK@0.2と採用率を同時に報告する．

\section{本研究の貢献}

本例で想定する貢献は，次の3点である．

\begin{enumerate}
  \item 関節信頼度と正規化骨格辺長を統合した，推定器非依存の骨格品質スコアを定義した．
  \item 動画系列単位で学習・検証・評価を分離し，低照度および動きぼけ条件に対する評価手順を示した．
  \item 除外によるPCK@0.2の変化だけでなく採用率も併記し，後段処理へ与える欠損とのトレードオフを整理した．
\end{enumerate}

これらは卒業論文の記述例として設定した架空の貢献であり，実在する研究成果ではない．

\section{本論文の構成}

第2章では人物姿勢推定と誤推定除外に関する技術的背景を述べ，第3章で入力，出力および評価課題を定義する．第4章では骨格品質スコアと判定処理を示し，第5章で架空データを用いた比較結果を報告する．第6章では得られた知見と限界をまとめ，今後の課題を示す．
